% Exposé I. The pointed slice
% Sources: theory/framework.md (Exposé I), site/sections/02-idea.html (#idea-methods ... #idea-dictionary),
% public statements of I.4, I.6, I.7, I.8 in theory/propositions.json. Proofs are in Appendix C.
\section{The pointed slice}\label{sec:slice}

We model a weight-space PEFT method as a pointed map into the pretrained weights. The frozen model and full fine-tuning sit
at the two ends of an interval, and sums and products of such maps already give methods in use. This exposé sets up the
category and its first invariants. The theorems begin when we look closely at one image near one point, which is where
Exposé~\ref{sec:image} starts.

% ------------------------------------------------------------------------------------------------
\subsubsection{Pretrained models and methods}\label{sec:slice-methods}

A pretrained network is a function with a parameter slot and a chosen value in it. Its natural home is the bicategory of
parametrised maps (\bgref{B.16}).

\begin{definition}{I.1}{Para; pretrained model}
For a cartesian category $\mathcal C$, the bicategory $\Para(\mathcal C)$ is defined as follows
\citep{fong2017backprop,cruttwell2021categorical}:
\begin{itemize}
  \item its 1-cells $X\to Y$ are pairs $(P,\,f:P\times X\to Y)$;
  \item its 2-cells $(P,f)\Rightarrow(Q,g)$ are maps $r:Q\to P$ with $g=f\circ(r\times X)$.
\end{itemize}
A \textbf{pretrained model} is a pointed 1-cell $F=(\Theta,\theta_0,f)$ of $\Para(\Diff)$. Its \textbf{realisation} is
$\Phi_f:\Theta\to C^\infty(X,Y)$, $\theta\mapsto f(\theta,-)$. Fine-tuning is a path $\theta(t)$ with $\theta(0)=\theta_0$.
\end{definition}

$\Para$ and the lens reading of backpropagation are due to \citet{fong2017backprop} and \citet{cruttwell2021categorical};
\citet{gavranovic2024position} develop a broader programme on the same foundations.

\begin{definition}{I.2}{methods and simulations}
A weight-space \textbf{PEFT method} at $\theta_0$ is a pointed smooth map $M=(Q,q_0,\rho)$, $\rho:Q\to\Theta$, with
$\rho(q_0)=\theta_0$. Usually $\dim Q\ll\dim\Theta$, and $\theta_0$ (or a splitting of it) is frozen inside $\rho$. These
form the pointed slice
\begin{equation*}
  \Meth(\Theta,\theta_0):=\Diff_*/(\Theta,\theta_0).
\end{equation*}
\begin{itemize}
  \item A \textbf{morphism} $h:M\to N$ is a pointed smooth map with $\rho_N\circ h=\rho_M$. Read it as ``$N$ simulates $M$''.
    Then $N$ reproduces every weight that $M$ reaches, together with its parameter path. A \emph{weak} morphism drops the
    condition $h(q_0)=q_0'$.
  \item The \textbf{budget} is $\lvert M\rvert=\dim Q$.
  \item $M$ is \textbf{additive} if $\rho=\theta_0+\delta$ with $\delta(q_0)=0$ and $\delta$ independent of $\theta_0$.
  \item $M$ is of \textbf{action type} if $\rho(q)=\gamma(q)\cdot\theta_0$ for a map $\gamma$ into a group acting on $\Theta$.
\end{itemize}
\end{definition}

\begin{example*}{four objects}
On one $m\times n$ box, LoRA \citep{hu2021lora} is $\rho(B,A)=W_0+sBA$ on $Q=\R^{m\times r}\times\R^{r\times n}$, pointed
at $(0,A_0)$. It is additive, with budget $r(m+n)$. OFT \citep{qiu2023controlling} with the exact Cayley map
(\bgref{D.7}) is $\rho(Q)=W_0\,C(Q)$ for skew $Q$ (block-diagonal in OFT), pointed at $Q=0$. It is of action type, with
$\gamma=C$ landing in $\SO(n)$ on the input side. The frozen model is $(\mathrm{pt},*,\theta_0)$, the inclusion of the
single point $\theta_0$, and full fine-tuning is $(\Theta,\theta_0,\id)$, the identity of $\Theta$.
\end{example*}

\begin{example*}{one arrow}
Take LoRA-FA \citep{zhang2023lora}, $B\mapsto W_0+sBA_0$ with $A_0$ frozen, pointed at $B=0$. The pointed map
$h(B)=(B,A_0)$ satisfies $\rho_{\rm LoRA}\circ h=\rho_{\rm FA}$, so LoRA pointed at $(0,A_0)$ simulates LoRA-FA
(Figure~\ref{fig:interval}c).
\end{example*}

Random-subspace training, $\rho(z)=\theta_0+Pz$ with a fixed random matrix $P$, is how \citet{li2018measuring} measured
intrinsic dimension and how \citet{aghajanyan2020intrinsic} fine-tuned language models in small subspaces. In the slice it
is a linear object, with $d(M)=\lvert M\rvert$ (\thmref{Definition}{I.5}) when $P$ has full column rank.

\begin{figure}[tbp]
\centering
{\small
\begin{tikzcd}[column sep=6.2em,row sep=0.5em]
  (\mathrm{pt},*) \arrow[r,"{*\,\mapsto\, q_0}"] \arrow[rr,bend left=24,inkgray,"{*\,\mapsto\,\theta_0}"]
  & (Q_M,q_0) \arrow[r,tide,"\rho_M","\textsf{merge}"']
  & (\Theta,\theta_0) \\
  |[inkgray]| \textsf{Frozen (initial)} & |[inkgray]| M & |[inkgray]| \textsf{Full FT (terminal)}
\end{tikzcd}}

\smallskip
{\sffamily\footnotesize\color{inkgray}(a) the interval [frozen, full FT]}

\bigskip
\begin{minipage}[b]{0.42\linewidth}\centering
{\small
\begin{tikzcd}[column sep=1.6em,row sep=3.0em]
  (Q_M,q_0) \arrow[rr,"h"] \arrow[dr,tide,"\rho_M"'] & & (Q_N,q_0') \arrow[dl,tide,"\rho_N"] \\
  & (\Theta,\theta_0) &
\end{tikzcd}}

\smallskip
{\sffamily\footnotesize\color{inkgray}(b) $N$ simulates $M$: $\rho_N\circ h=\rho_M$}
\end{minipage}\hfill
\begin{minipage}[b]{0.56\linewidth}\centering
{\small
\begin{tikzcd}[column sep=0.6em,row sep=3.0em]
  (\R^{m\times r},0) \arrow[rr,"{B\,\mapsto\,(B,A_0)}"] \arrow[dr,tide,"{W_0+sBA_0}"']
  & & (\R^{m\times r}\times\R^{r\times n},(0,A_0)) \arrow[dl,tide,"{W_0+sBA}"] \\
  & (\R^{m\times n},W_0) &
\end{tikzcd}}

\smallskip
{\sffamily\footnotesize\color{inkgray}(c) LoRA simulates LoRA-FA}
\end{minipage}
\caption{The pointed slice $\Meth(\Theta,\theta_0)$.
(a)~The interval of weight-space methods. For a method $M$, the first arrow picks the initialisation $q_0$; it lies over
$\Theta$ because $M$ starts at $\theta_0$. The second arrow is $\rho_M$ itself, the merge into the base weights, and the
composite is the inclusion of $\theta_0$. The frozen model is initial and full fine-tuning terminal
(\thmref{Observation}{I.4}). A pointing defect, $\rho_M(q_0)\neq\theta_0$, leaves the triangle commuting only up to that
defect (Exposé~\ref{sec:base}). Adapters and prefixes fall outside the interval because they extend the network
(\thmref{Definition}{VI.1}).
(b)~A simulation is a pointed smooth map $h$ that makes the triangle over $(\Theta,\theta_0)$ commute; a weak simulation
drops $h(q_0)=q_0'$.
(c)~An instance: LoRA pointed at $(0,A_0)$ simulates LoRA-FA through $h(B)=(B,A_0)$.}
\label{fig:interval}
\end{figure}

% ------------------------------------------------------------------------------------------------
\subsubsection{What Para adds}\label{sec:slice-para}

\begin{remark}{I.3}{orientation, and what Para adds}
Each $\rho$ is a 2-cell $F\Rightarrow F^\rho:=(Q,f\circ(\rho\times X))$ \emph{out of} $F$. A simulation $h:M\to N$ is a
2-cell $F^{\rho_N}\Rightarrow F^{\rho_M}$, since $\rho_N\circ h=\rho_M$. Hence $\Meth(\Theta,\theta_0)$ is the
\textbf{opposite of the pointed coslice of reparametrisation 2-cells under $F$} (\bgref{A.7}). The slice over $F$ is a
different category, and its objects are extensions (Exposé~\ref{sec:merge}).

$\Meth$ depends only on $(\Theta,\theta_0)$ and not on $f$. So for weight-space methods, $\Para$ is bookkeeping. The
network enters in exactly three places, and we use $\Para$ only there:
\begin{itemize}
  \item extensions and mergeability (Exposé~\ref{sec:merge}; \thmref{Definition}{VI.1}, \thmref{Theorem}{VI.3});
  \item the realisation $\Phi_f$, through functional images and descent (\thmref{Observation}{II.15});
  \item the backpropagation lens (Exposé~\ref{sec:lens}; \thmref{Definition}{IV.1}, \thmref{Proposition}{IV.2}).
\end{itemize}
\end{remark}

Calling $\Para$ bookkeeping concerns this slice only. In the three places listed we use the framework of
\citet{fong2017backprop} and \citet{cruttwell2021categorical} directly, including their lens 2-functor in
Exposé~\ref{sec:lens}.

This is also why the next exposés work in weight space. Reach in $\Theta$ is a property of the image, and the cometric
$K=D\rho\,\Lambda\,D\rho\T$ that shapes training is built from $\rho$ and the optimizer's rates $\Lambda$. The network
enters training through the gradient it hands to the method, and enters expressivity through $\Phi_f$ once we ask about
functions instead of weights. So Exposés~\ref{sec:image} and~\ref{sec:fibres} work in $\Theta$ and flag the places where
$f$ matters.

% ------------------------------------------------------------------------------------------------
\subsubsection{Frozen and full fine-tuning}\label{sec:slice-interval}

\begin{observation}{I.4}{universal properties}
In $\Meth(\Theta,\theta_0)$:

(a) the frozen model $(\mathrm{pt},*,\theta_0)$ is initial, and full fine-tuning $(\Theta,\theta_0,\id)$ is terminal. The
unique morphism $M\to{}$Full FT is $\rho_M$ itself.

(b) If the fibre product $Q_M\times_\Theta Q_N$ (\bgref{A.11}) exists in $\Diff$, it is the product $M\times N$. It exists,
for instance, when the set-level fibre product is an embedded submanifold of $Q_M\times Q_N$: a clean intersection
(\bgref{C.7}), or linear or affine $\rho$'s as in \thmref{Proposition}{I.8}(a). Transversality $\rho_M\pitchfork\rho_N$ is
never the reason when $\lvert M\rvert+\lvert N\rvert<\dim\Theta$, since at $(q_0,q_0')$ it would force
$\lvert M\rvert+\lvert N\rvert\ge\dim\Theta$. In $\Set$ the fibre product always exists and
$\Im(M\times N)=\Im M\cap\Im N$. Whenever the product exists in $\Meth$, its underlying set is this set-level fibre
product, so the image identity holds in every case. The product can fail to exist. On $\Theta=\R$ ($1\times1$ matrices),
$\mathrm{LoRA}_1\times\text{Frozen}$ would have the cross $\{BA=0\}\subset\R^2$ as underlying set. No manifold structure
on the cross has the universal property, because both axes would have to lift smoothly through the point over the origin,
which forces rank $2$ there.

(c) $\Im$ is a functor (\bgref{A.4}) to the poset (\bgref{A.16}) of subsets of $\Theta$ containing $\theta_0$. Conversely,
$\Im M\subseteq\Im N$ yields a morphism of the set-level slice $\Set_*/(\Theta,\theta_0)$, though not necessarily a smooth
one. Take $\rho_M(t)=t$ and $\rho_N(u)=u^3$ on $\R$, pointed at $0$. The images agree, yet no smooth $h:\R\to\R$ satisfies
$\rho_N\circ h=\id$, that is, $h(t)^3=t$.
\end{observation}

\noindent The proof is in Appendix~\ref{proof:I.4}. The obstruction in (b) covers every pointing of $\mathrm{LoRA}_1$,
including the standard one $(0,a_0)$ with $a_0\neq0$, whose base point is a smooth point of the cross; the obstruction sits
at the origin.

\begin{explains}
Every weight-space method pointed at $\theta_0$ lies in the interval [frozen, full FT]. Adapters, prefixes and ReFT
\citep{wu2024reft} are extensions outside $\Meth$ (Exposé~\ref{sec:merge}). ``Merge into the base weights'' is the
canonical morphism to the terminal object. In exact arithmetic on a full-precision base this is definitional, since every
object of $\Meth$ is mergeable by construction. Type-preserving merges into quantised bases are lossy
(\thmref{Proposition}{V.2}(d)), combining several adapters is a separate question (\thmref{Proposition}{V.4}), and the
remaining zero-overhead questions live on extensions (Exposé~\ref{sec:merge}). What both $M$ and $N$ can express,
$\Im M\cap\Im N$, always exists as a set. The categorical product, when it exists, is the universal method simulated by
both, and its image is exactly this set. \thmref{Proposition}{I.8}(a) computes a genuine one (LoRA-XS).
\thmref{Proposition}{I.8}(b) and \thmref{Theorem}{II.7} compute meets, that is, intersections in the poset of subsets of
$\Theta$ containing $\theta_0$; they are not products in $\Meth$.
\end{explains}

\paragraph{Defaults that do not start at the checkpoint}
Some defaults move the weights before training starts. In Hugging Face PEFT \citep{xmangrulkar2022peft} the default
FourierFT \citep{gao2024parameter} draws its spectrum from a standard normal, and the default vector bank and logits of
VB-LoRA \citep{li2024vb} are random too. LoRA with \texttt{init\_lora\_weights=False}, LoRA-XS with its small random $R$, and every method on
a quantised base also start away from $\theta_0$. Each lives over its own starting weight or carries a pointing defect,
which Exposé~\ref{sec:base} measures.

\paragraph{Lifting obstructions beyond first order}
The gap between $\Set$ and $\Diff$ in (c) is a smooth lifting obstruction. \thmref{Observation}{I.6} below detects its
first-order instances. In the cube-root example, $S_1(M)=\R\not\subseteq S_1(N)=0$. Higher-order obstructions need finer
invariants. The maps $\rho_M(t)=t^3$ and $\rho_N(u)=u^5$ have equal images and $S_1=0$ on both sides, yet the only lift,
$t^{3/5}$, is not $C^1$.

% ------------------------------------------------------------------------------------------------
\subsubsection{Invariants of a method}\label{sec:slice-invariants}

The budget $\lvert M\rvert$ counts parameters. Two other numbers say more. One is the dimension of what a method
reaches; the other is the dimension of what it reaches to first order at the start, which is what the first gradient steps
use. For LoRA with $r=16$
on a $4096\times4096$ matrix the budget is $131{,}072$, the image has dimension $r(m+n-r)=130{,}816$, and from zero
initialisation the first steps use only $mr=65{,}536$ directions (\thmref{Theorem}{II.2}, \thmref{Proposition}{II.12}).

\begin{definition}{I.5}{invariants of an object}
For $M=(Q,q_0,\rho)$:
\begin{itemize}
  \item the image $\Im M=\rho(Q)$;
  \item the \textbf{expressive dimension} $d(M)=\dim\Im M=\max_q\rk d\rho_q$;
  \item the \textbf{gauge waste} $g(M)=\lvert M\rvert-d(M)$;
  \item the \textbf{first-order image} $S_1(M)=d\rho_{q_0}(T_{q_0}Q)$, and the \textbf{first-order deficit}, which is
    $d(M)-\dim S_1(M)$;
  \item the \textbf{gauge group} $\{\varphi\in\mathrm{Diff}(Q):\rho\circ\varphi=\rho\}$, which is unpointed, since a gauge
    transformation moves $q_0$ inside its fibre. In practice we exhibit a Lie group (\bgref{D.1}) acting on $Q$ that
    preserves $\rho$ and is locally transitive on generic fibres, and call it the \textbf{generic gauge}.
\end{itemize}
The base point $\theta_0$ is \textbf{regular} for $M$ if the deficit is $0$, and an \textbf{apex} if the deficit is
positive. We say \textbf{vertex} for the subset-level notion, the cone point of an image germ (\bgref{C.10}), and keep
\textbf{apex} for the base-point notion. The two differ in general. At a crossing of two smooth branches the image is
singular, yet the deficit can be $0$. In every case below with a positive deficit, the image germ is (diffeomorphic to) a
cone with vertex $\theta_0$; the exception is DoRA \citep{liu2024dora}, which inherits LoRA's deficit through its torus
factor. A cone germ that is a $C^1$ manifold at its vertex equals its tangent space (\bgref{C.2}) there, so it is linear.
Hence the vertex of a non-linear cone is a singular point.
\end{definition}

\begin{remark*}{the word ``gauge''}
``Gauge'' here means the symmetry group of a parametrisation, as in \thmref{Definition}{I.5}. There is no connection, and
away from the full-rank stratum the fibres jump. The word borrows, as an analogy, the physicist's picture of redundant
coordinates that leave every observation unchanged. Only the group and its action enter the theorems.
\end{remark*}

If $N$ simulates $M$, whatever $M$ can do at its first step, $N$ can do too.

\begin{observation}{I.6}{$S_1$ is functorial}
If $M\to N$ is a (pointed) morphism, then $S_1(M)\subseteq S_1(N)$. In particular, isomorphic methods have equal $S_1$ and
equal first-order deficit. Monotonicity holds only along pointed morphisms. Along a weak morphism, such as a gauge map
moving $q_0$ inside its fibre, $S_1$ can change.
\end{observation}

\noindent The proof is in Appendix~\ref{proof:I.6}.

\begin{explains}
$S_1$ is the set of weight-space directions available to the linearised (lazy, NTK) dynamics at initialisation, since the
lazy-regime kernel $J_f\,d\rho\,\Lambda\,d\rho\T J_f\T$ has range $J_f(S_1)$. That kernel also depends on $\Lambda$ and on
the network Jacobian $J_f$, not only on $S_1$. For random-$A_0$ zero-init LoRA it approximates the full fine-tuning kernel
\citep{malladi2022kernel}, so a large first-order deficit need not cost much on the data. $S_1$ is an isomorphism
invariant that separates initialisations the image cannot, namely zero-init $\mathrm T_A$ against $\mathrm T_B$, and the
choice of $\operatorname{row}A_0$ (\thmref{Theorem}{II.2}). The image already separates split initialisations, since it
determines $P$.
\end{explains}

The kernel view of language-model fine-tuning, and the closeness of LoRA's kernel to full fine-tuning's, are due to
\citet{malladi2022kernel}. \citet{hayou2024impact} compared the zero initialisations $B_0=0$ and $A_0=0$ as training
procedures; in the slice they are already different objects.

% ------------------------------------------------------------------------------------------------
\subsubsection{Sums}\label{sec:slice-sums}

Weight space is a vector space, so running two methods side by side and adding their updates gives a method again.

\begin{observation}{I.7}{additive monoidal structure}
$\Theta$ is a vector space, so every object of $\Meth(\Theta,\theta_0)$ can be written $\rho=\theta_0+\delta$ with
$\delta:=\rho-\theta_0$ and $\delta(q_0)=0$. Put
\begin{equation*}
  M\oplus N:=\big(Q_M\times Q_N,\ (q_0,q_0'),\ \theta_0+\delta_M+\delta_N\big).
\end{equation*}
This is a symmetric monoidal structure (\bgref{B.1}) with the frozen model as unit. Its image is the Minkowski sum
$\theta_0+(\Im\delta_M+\Im\delta_N)$ (\bgref{B.14}). It differs from the categorical product, whose image, when the
product exists, is the intersection $\Im M\cap\Im N$ (\thmref{Observation}{I.4}(b)). ``Additive'' in
\thmref{Definition}{I.2} ($\delta$ independent of $\theta_0$) is a condition on \emph{families}
$\theta_0\mapsto M(\theta_0)$; it matters for base change (Exposé~\ref{sec:base}) and plays no role here.

Write $\mathrm{LoRA}_k$ for LoRA with frozen residual $\theta_0$; for split pointings the sum below has residual
$\theta_0-P-P'$. Concatenation $h(B_1,A_1,B_2,A_2)=([B_1\ B_2],[A_1;A_2])$ is an isomorphism
$\mathrm{LoRA}_r\oplus\mathrm{LoRA}_s\cong\mathrm{LoRA}_{r+s}$ of unpointed objects. It is a pointed isomorphism onto
$\mathrm{LoRA}_{r+s}$ pointed at $([B_0\ B_0'],[A_0;A_0'])$. For $\mathrm T_A$ pointings ($B_0=0$, $\rk A_0=r$), the
target pointing has $\rk[A_0;A_0']=r+s$ if and only if $\operatorname{row}A_0\cap\operatorname{row}A_0'=0$; it lies in the
stratum $\mathrm T_A$ of \thmref{Theorem}{II.2} when moreover $r+s<\min(m,n)$.

Standard pointings ($\mathrm T_A$, or symmetrically $\mathrm T_B$) exist only for $r<\min(m,n)$ (\thmref{Theorem}{II.2}),
so the monoidal question is pointwise. If $F(r)$ and $F(2r)$ are standard ($2r<\min(m,n)$), then
$F(r)\oplus F(r)\not\cong F(2r)$ as pointed objects; for $2r\ge\min(m,n)$ no standard $F(2r)$ exists. Hence no assignment
$r\mapsto\mathrm{LoRA}_r$ with $F(1)$ and $F(2)$ standard is \textbf{strong monoidal} (\bgref{B.15}) into pointed objects.
For $r\neq s$ and standard pointings, $F(r)\oplus F(s)\cong F(r+s)$ as pointed objects exactly when the row space of the
target's $A_0$ is the direct sum of the two source row spaces. The assignment is strong monoidal on unpointed objects, up to
weak morphisms, and on pointed objects with degenerate pointings, namely the origin $(B_0,A_0)=(0,0)$, where $S_1=0$, or
the stacked pointings $F(r)=F(1)^{\oplus r}$. Concatenation is strictly associative and unital, so discrete $(\N,+)$
suffices for a plain strong monoidal functor. \emph{Symmetric} strong monoidality needs the permutation groupoid, since
swapping and then concatenating permutes the blocks.

With per-adapter scales $s_r$ ($\alpha/r$ for LoRA, $\alpha/\sqrt r$ for rsLoRA) and any target scale $s_*$, the map
$h=([B_1\ B_2],[c_1A_1;c_2A_2])$ with $c_i=s_{r_i}/s_*$, where $r_1=r$ and $r_2=s$, is an isomorphism over $\Theta$. It is
a linear bijection and preserves row spaces, so the pointing criterion is unchanged.
\end{observation}

\noindent The proof is in Appendix~\ref{proof:I.7}. The negative claim rests on $\dim S_1$, which any pointed
isomorphism preserves (\thmref{Observation}{I.6}). It is $mr$ for $F(r)\oplus F(r)$ and $2mr$ for a standard $F(2r)$.
[Num] At $m=7$, $\dim S_1=28$ for a standard $\mathrm{LoRA}_4$, against $14$ for $\mathrm{LoRA}_2\oplus\mathrm{LoRA}_2$
with the same $A_0$ and $28$ with independent $A_0$'s; $F(1)\oplus F(2)$ and a standard $F(3)$ both give $21$
(Appendix~\ref{app:repro}).

\paragraph{The \texttt{cat} merge of Hugging Face PEFT}
Hugging Face PEFT's \texttt{add\_weighted\_adapter(..., combination\_type="cat")} \citep{xmangrulkar2022peft} evaluates the
rescaled concatenation $h$ at trained points. It takes $s_*=1$. It sets the merged \texttt{lora\_alpha} to the new rank
$r+s$, multiplies each $A$ block by $w_is_{r_i}$, so that $c_i=w_is_{r_i}$, and leaves the $B$ blocks alone. It is the
isomorphism $h$ when all weights $w_i=1$ and the merged adapter has \texttt{use\_rslora=False}. In PEFT releases up to
v0.21.1 the merged adapter inherits \texttt{use\_rslora} from the first source. For rsLoRA sources
\citep{kalajdzievski2023rank} its scale is then $(r+s)/\sqrt{r+s}=\sqrt{r+s}$, and the merged update is $\sqrt{r+s}$ times
too large. PR \#3449 on the main branch fixes this by setting \texttt{use\_rslora=False} on the merged adapter.

\begin{explains}
RoSA \citep{nikdan2024rosa} (LoRA $\oplus$ a fixed-mask sparse method) and GraLoRA \citep{jung2025gralora} ($\oplus$ of
$k^2$ block-embedded LoRAs) are $\oplus$-words, and their expressivity is a Minkowski sum. ReLoRA \citep{lialin2023relora}
is $\mathrm{LoRA}^{\oplus K}$ unrolled in time, apart from its full-rank warm start. A sum of $k$ rank-$r$ LoRAs, each with
residual $\theta_0$, is isomorphic in $\Meth$ to one LoRA of rank $kr$, pointed at the concatenated initialisation. Under
$\alpha/r$ scaling the isomorphism rescales the $A$ blocks, so training at equal $\alpha$ and learning rate follows a
different trajectory. The \texttt{cat} merge is the isomorphism $h$ evaluated at trained points; the merged adapter, if
unfrozen, is $\mathrm{LoRA}_{r+s}$.
\end{explains}

\paragraph{What is and is not a sum}
LoRA-soup CAT \citep{prabhakar2024lora} freezes the trained modules and learns layer-wise weights, so it is a $\oplus$ of
$k$ frozen line methods $\alpha_i\mapsto\alpha_iB_iA_i$. Its per-layer image is a span of dimension at most $k$
(generically exactly $k$), much smaller than the image of $\mathrm{LoRA}_{kr}$. UniPELT \citep{mao2021unipelt} and MAM
\citep{he2021unified} are \textbf{not} $\oplus$-words in $\Meth(\Theta,\theta_0)$. Their adapter and prefix parts live on
extensions (Exposé~\ref{sec:merge}), and UniPELT's gates depend on the input. The obstruction in
\thmref{Observation}{I.7} is the reuse of one pointing. With a fixed pointing, stacking the same $A_0$ twice gives a
rank-deficient frame, which $S_1$ distinguishes from a standard $\mathrm{LoRA}_{2r}$.

% ------------------------------------------------------------------------------------------------
\subsubsection{Products}\label{sec:slice-products}

Freeze $A$ at $A_0$ and train $B$, or freeze $B$ at $B_0$ and train $A$. The product of these two one-sided methods is
LoRA-XS. As a pullback square over $\Theta$ it reads
\begin{equation*}
\begin{tikzcd}[column sep=6.5em,row sep=2.8em]
  \R^{r\times r} \arrow[r,"{R\,\mapsto\, B_0R}"] \arrow[d,"{R\,\mapsto\, RA_0}"'] \arrow[dr,phantom,"\lrcorner",pos=0.09]
  & Q_{\rm FA}=\R^{m\times r} \arrow[d,tide,"{B\,\mapsto\, W_0+BA_0}"] \\
  Q_{\rm FB}=\R^{r\times n} \arrow[r,tide,"{A\,\mapsto\, W_0+B_0A}"] & \Theta=\R^{m\times n}
\end{tikzcd}
\end{equation*}
and both composites send $R$ to $W_0+B_0RA_0$.

\begin{proposition}{I.8}{a product that produces a named method; an image-level meet}
(a) Let $A_0\in\R^{r\times n}$ have full row rank and $B_0\in\R^{m\times r}$ full column rank. Define the one-sided methods
\begin{equation*}
  \mathrm{FA}(A_0):B\mapsto W_0+BA_0,\qquad \mathrm{FB}(B_0):A\mapsto W_0+B_0A,
\end{equation*}
both pointed at $0$. Then
\begin{equation*}
  \mathrm{FA}(A_0)\times\mathrm{FB}(B_0)\;\cong\;\big(\R^{r\times r},\ 0,\ R\mapsto W_0+B_0RA_0\big).
\end{equation*}
The product supplies this sandwich for \emph{any} frames. Assume $\rk W_0\ge r$ ($\sigma_r>0$) and take the SVD frames
$B_0=U_r$, $A_0=\Sigma_rV_r\T$. The result is \textbf{LoRA-XS pointed at $R=0$}.

(b) Let $W_0$ have full row rank (so $m\le n$), and let $\mathrm{LoRA}_r$ carry its standard pointing ($B_0=0$, image
$W_0+\Mr$). The image-level meet of $\mathrm{LoRA}_r$ with the output-scaling method $\ell\mapsto\diag(\ell)W_0$ is
$\{\diag(\ell)W_0:\#\{i:\ell_i\neq1\}\le r\}$, that is, ``rescale at most $r$ output neurons''.

\upshape The LoRA-XS paper \citep{baazy2024lora} writes $xW+xARB$ in row-vector convention with $A=U_r\Sigma_r$ and
$B=V_r\T$. Transposed, the update is $B_0R\T A_0$, so the isomorphism is $R\mapsto\alpha R\T$. LoRA-XS initialises
$R\sim\mathcal N(0,\sigma^2)$ with $\sigma=10^{-5}$, a negligible pointing defect. The output-scaling method in (b) is the
key/value part of (IA)$^3$ \citep{liu2022few}. On the FFN, (IA)$^3$ scales the \emph{input} of $W_{\rm down}$,
$W\mapsto W\diag(\ell)$; the mirror statement then needs full column rank, or holds generically for $r<d$, where $d$ is the
number of rows of $W_{\rm down}$.
\end{proposition}

\noindent The proof is in Appendix~\ref{proof:I.8}. In (a) the fibre product is the $r^2$-dimensional linear subspace
$\{(B_0R,RA_0)\}$, and if $\rk W_0<r$ the frame $\Sigma_rV_r\T$ is rank-deficient and the LoRA-XS instance breaks.
[Num] The fibre product of (a) has dimension $4=r^2$ at $(m,n,r)=(6,8,2)$ (Appendix~\ref{app:repro}).

\begin{explains}
LoRA-XS is the largest method that factors through \emph{both} frozen frames. In the frames of an initialisation it is the
overlap that a split initialisation double-counts (\thmref{Proposition}{II.3}); it is LoRA-XS proper when those frames span
the top-$r$ singular subspaces of $W_0$, as at a PiSSA initialisation \citep{meng2024pissa}. The meet in (b) is only
set-level. For $n\ge2$ the product does not exist in $\Diff$, because its fibre product is singular at the origin
$(b,a,\ell)=(0,0,\mathbf 1)$, by the rank argument of \thmref{Observation}{I.4}(b) (Appendix~\ref{proof:I.8}). Once a support is fixed, the meet is
just (IA)$^3$ masked to $r$ neurons; compare sparse activation scaling \citep{xstoehr2024activation}.
\end{explains}

\paragraph{The shape of the meet in (b)}
The meet is the union of the $\binom mr$ coordinate $r$-planes ``rescale at most $r$ output neurons''. For $r=1$ the fibre
product is the union of the linear spaces $\{b=0,\ \ell=\mathbf 1\}$ and $\{a=0,\ \ell=\mathbf 1\}$ with $m$
two-dimensional sheets $\{b=ce_i,\ a=\lambda w_i,\ \ell_i=1+sc\lambda\}$, where $w_i$ is row $i$ of $W_0$. Under the
standard pointing ($b=0$, $a_0$ generic) the base point is a smooth point of $\{b=0,\ \ell=\mathbf 1\}$, so the singularity
sits away from it, at the origin.

% ------------------------------------------------------------------------------------------------
\subsubsection{The dictionary so far}\label{sec:slice-dictionary}

Table~\ref{tab:slice-dictionary} collects the translations made in this exposé. The full dictionary is the Rosetta stone
of Exposé~\ref{sec:rosetta}.

\begin{table}[!htb]
\centering
\caption{The dictionary for Exposé~\ref{sec:slice}.}
\label{tab:slice-dictionary}
\small\renewcommand{\arraystretch}{1.15}
\begin{tabularx}{\linewidth}{>{\raggedright\arraybackslash}p{0.27\linewidth}>{\raggedright\arraybackslash}X>{\raggedright\arraybackslash}p{0.2\linewidth}}
\toprule
\headrow \textbf{In PEFT} & \textbf{In the pointed slice} & \textbf{Statement} \\
\midrule
pretrained checkpoint & base point $\theta_0$ of a pointed 1-cell $(\Theta,\theta_0,f)$ & \thmref{Definition}{I.1} \\
a weight-space PEFT method & an object $\rho:(Q,q_0)\to(\Theta,\theta_0)$ of $\Meth(\Theta,\theta_0)$ & \thmref{Definition}{I.2} \\
``$N$ can do whatever $M$ can'' & a morphism $h:M\to N$, $\rho_N\circ h=\rho_M$; at the level of sets, image inclusion
  & \thmref{Definition}{I.2}, \thmref{Observation}{I.4}(c) \\
frozen model, full fine-tuning & initial and terminal objects & \thmref{Observation}{I.4}(a) \\
merging into full-precision weights & the canonical morphism $\rho_M$ to the terminal object & \thmref{Observation}{I.4}(a) \\
effective number of parameters & expressive dimension $d(M)$; the rest is gauge waste & \thmref{Definition}{I.5} \\
what the first steps can use & first-order image $S_1$, a functorial invariant & \thmref{Observation}{I.6} \\
RoSA, GraLoRA, ReLoRA stages & $\oplus$-words, with a Minkowski sum as image & \thmref{Observation}{I.7} \\
HF \texttt{cat} merge & the concatenation isomorphism evaluated at trained points & \thmref{Observation}{I.7} \\
LoRA-XS & the product $\mathrm{FA}(A_0)\times\mathrm{FB}(B_0)$ & \thmref{Proposition}{I.8} \\
adapters, prefixes, ReFT & extensions of the network, outside $\Meth$ & \thmref{Definition}{VI.1} \\
\bottomrule
\end{tabularx}
\end{table}
